\documentclass[10pt,a4paper]{amsart}
\usepackage[margin=19mm]{geometry}
\usepackage{amsmath,amssymb,mathtools}
\usepackage[scr=rsfs,cal=euler]{mathalfa}
\usepackage{xfrac}
\usepackage[unicode,hidelinks,bookmarks=false]{hyperref}
\newcommand{\ie}{i.e.\ }
\newcommand{\cf}{cf.\ }
\newcommand{\resp}{respectively\ }
\newcommand{\Subsection}{\subsection}
\providecommand{\nobreakdash}{\nobreak\hbox{-}}
\setlength{\emergencystretch}{2em}
\multlinegap0pt
\usepackage[shortlabels]{enumitem}
\usepackage{mathtools}
\usepackage{amssymb}
\usepackage{xcolor}
\usepackage{bbm}
\usepackage{bm}
\usepackage{tikz}
\usepackage{float}
\usepackage[all]{xy}
\newtheorem{lemma}{Lemma}
\def\R{\mathbb{R}}
\title{Chow polynomials of rank-uniform labeled posets}
\author{Basile Coron, Luis Ferroni, Shiyue Li}
\date{Source: arXiv:2511.13819v2 (CC BY 4.0)}
\begin{document}
\maketitle

\noindent\emph{Source and licence.} Chow polynomials of rank-uniform labeled posets. Version arXiv:2511.13819v2 (CC BY 4.0), distributed by arXiv under the Creative Commons Attribution 4.0 International licence, http://creativecommons.org/licenses/by/4.0/, as recorded in the arXiv licence metadata for this exact version and shown on the abstract page https://arxiv.org/abs/2511.13819v2. Full licence notice: \texttt{licenses/arxiv-2511.13819-CC-BY-4.0.txt}.\par\smallskip

\noindent\textit{Editorial scope.} Counted unit: the article's lemma lemma (source0.tex lines 702--719). The article's complete original proof of it is delivered with it (source0.tex lines 721--791, one \texttt{proof} environment of 71 lines). No mathematics is written, repaired or re-derived: the statement and the proof are the article's own lines, in the article's own notation, and no retained line is substituted or re-flowed.  No further source-local context is needed: every cross-reference of every retained line resolves inside the two delivered spans.  Nothing was omitted from any retained span, so the package carries no omission record.  No retained line contains a citation, so the wrapper carries no bibliography and no bibliography entry.  No retained line contains a figure environment, an image-inclusion command or drawing code, so no image is delivered and the package declares no asset dependency.  Editorial formatting only, in addition: an AMS \texttt{amsart} preamble that restates the article's own declarations and macros used by the retained lines, namely the environments \texttt{lemma} on separate counters, together with the standard packages those lines need.  The retained environments print in this document as Lemma 1 in order of appearance, whereas the article prints the same items with the numbering of its own full document.  The complete native source of the article as distributed in the arXiv e-print for this exact version is bundled unchanged as \texttt{sources/189319/source0.tex}.
\begin{lemma}\label[lemma]{lem:two-column-monotonicity-interlacing}
    Let $m$ and $n$ be positive integers, and let $d_1 \leqslant d_2 \leqslant \cdots \leqslant d_{m}$ be a weakly increasing sequence of integers. 
    Let $\mathbf{A}(x)$ be an $m \times n$ matrix with polynomial entries defined as 
    \[
    \mathbf{A}(x) \coloneqq \begin{pmatrix} a_{ij}(x) \end{pmatrix}, \text{ where } a_{ij}(x) = \begin{cases}
    \alpha_{i} \cdot x & \text{ if } j \leqslant d_i, \\
    \alpha_{i} & \text{ if } j > d_i, 
    \end{cases} \quad \alpha_i \in \R_{\geqslant 0}.
    \] 
    For any fixed column number $1 \leqslant j \leqslant n$, any two rows $1 \leqslant i < i' \leqslant m$, and nonnegative real numbers $p, q$, we write $f_j(x)$ for the polynomial 
    \begin{align*}
        f_{j}(\mathbf{A}; p, q, i, i'; x) \coloneqq (p \: x + q ) \:  a_{i \: j}(x) + a_{i' \: j}(x). 
    \end{align*}
    Then for any two column numbers $1 \leqslant j < j' \leqslant n$, and any nonnegative real numbers $p, q \in \R_{\geqslant 0}$,  we have that  
    \[
    f_{j'}\preceq f_{j}.
    \]
\end{lemma}

\begin{proof}
    By monotonicity of the sequence $(d_k)_{k=1}^{m}$, we have $d_i \leqslant d_{i'}$. 
    By definition of the matrix $A$, there are $5$ cases, depending on the relative positions of $j < j'$ with respect to $d_i \leqslant d_{i'}$. For each of the cases, we fix $\alpha_{i}, \alpha_{i'} \in \R_{\geqslant 0}$, and consider the $2 \times 2$ submatrix 
    \[
    \mathbf{A}_{i, \:i',\: j,\:j'} \coloneqq 
    \begin{pmatrix}
        a_{i\: j} & a_{i \: j'} \\
        a_{i' \: j} & a_{i' \: j'}
    \end{pmatrix}. 
    \]
    \begin{enumerate}[\normalfont(i), leftmargin=18pt]
        \item If $j < j' \leqslant d_i \leqslant d_{i'}$, then we have 
        \[
        \mathbf{A}_{i, \:i',\: j,\:j'}(x) = 
        \begin{pmatrix}
        \alpha_i \: x & \alpha_i \: x  \\
        \alpha_{i'} \: x & \alpha_{i'} \: x
        \end{pmatrix}.
        \]
        This implies that $f_{j} = f_{j'}$ for all $p, q \in \R_{> 0}$. 
        \item If $j \leqslant  d_i < j' \leqslant d_{i'}$, then we have 
        \[
        \mathbf{A}_{i, \:i',\: j,\:j'}(x) = 
        \begin{pmatrix}
        \alpha_i \: x & \alpha_i  \\
        \alpha_{i'} \: x & \alpha_{i'} \: x
        \end{pmatrix}.
        \]
        This implies that 
        \begin{align*}
            f_{j}(x) &= (px + q) \alpha_ix + \alpha_{i'} x, \text{ with roots } - \frac{\alpha_{i}q + \alpha_{i'}}{\alpha_{i}p}, \text{ and }0,\\
            f_{j'}(x) &= (px + q) \alpha_i + \alpha_{i'} x, \text{ with root } - \frac{\alpha_{i} q}{\alpha_{i}p + \alpha_{i'}}.\\
        \end{align*}
        This implies that $f_{j'} \preceq f_{j}$ for all $p, q \in \R_{> 0}$. 

        \item If $d_{i} < j < j' \leqslant d_{i'}$, then we have
        \[
        \mathbf{A}_{i, \:i',\: j,\:j'}(x) = 
        \begin{pmatrix}
        \alpha_i  & \alpha_i  \\
        \alpha_{i'} \: x & \alpha_{i'} \: x
        \end{pmatrix}.
        \]
        This implies that $f_{j'}(x) = f_{j}(x)$ for all $p, q \in \R_{> 0}$. 

        \item If $d_i < j \leqslant d_{i'} < j'$, then we have  
        \[
        \mathbf{A}_{i, \:i',\: j,\:j'}(x) = 
        \begin{pmatrix}
        \alpha_i  & \alpha_i  \\
        \alpha_{i'} \: x & \alpha_{i'}
        \end{pmatrix}.
        \]
        This implies that 
        \begin{align*}
            f_{j}(x) &= (px+q) \alpha_{i} + \alpha_{i'} x, \text{ with root} - \frac{\alpha_{i} q}{\alpha_{i} p + \alpha_{i'}},\\
            f_{j'}(x) &= (px+q) \alpha_{i} + \alpha_{i'}, \text{ with root} -\frac{\alpha_{i}q + \alpha_{i'}}{\alpha_{i}p}.  
        \end{align*}
        This implies that $f_{j'} \preceq f_{j}$ for all $p, q \in \R_{> 0}$.

        \item If $d_i \leqslant d_{i'} < j < j'$, then we have 
        \[
        \mathbf{A}_{i, \:i',\: j,\:j'} =
        \begin{pmatrix}
        \alpha_i  & \alpha_i  \\
        \alpha_{i'} & \alpha_{i'}
        \end{pmatrix}.
        \]
        This implies that $f_{j'} = f_{j}$ for all $p, q \in \R_{> 0}$. \qedhere
    \end{enumerate}
\end{proof}

\end{document}
